\documentclass[11pt]{article}
\usepackage{amsmath,amssymb,amsthm}
\newtheorem{theorem}{Theorem}
\newcommand{\ket}[1]{\lvert #1\rangle}
\begin{document}
\title{A four-band radical construction of second-order anticoherent subspaces}
\author{}\date{}\maketitle

This construction gives a general lower bound.  No assertion of optimality
for all $N$ is made.  The proof contains a finite, exactly reproducible
polynomial coefficient certificate, covering all $N$ simultaneously.

\begin{theorem}
For every integer $N\ge40$,
\[
 K(N,2)\ge\left\lfloor\frac N8\right\rfloor-3.
\]
A basis can be obtained using rational arithmetic and square roots.
\end{theorem}

Write $N=8q+\rho$, with $0\le\rho\le7$, and set
\[
 k=q-3,\quad h=2q,\quad s=4+\lfloor\rho/2\rfloor,\quad
 j=N/2,\quad v=\frac{j(j+1)}{3h^2},\quad x_i=\frac{j-s-2i}{h}.
\]
Define four rational weights for each $0\le i<k$ by evaluating
\begin{align*}
 w_0(x)&=\frac v4+\frac{x^2}4-\frac{35x}{36}+\frac56,&
 w_1(x)&=-\frac v4-\frac{x^2}4+\frac{5x}{12}+\frac12,\\
 w_2(x)&=-\frac v4-\frac{x^2}4+\frac{13x}{12}-\frac12,&
 w_3(x)&=\frac v4+\frac{x^2}4-\frac{19x}{36}+\frac16.
\end{align*}
The basis has the form
\[
 \ket{u_i}=\sum_{\ell=0}^3
 \sqrt{w_\ell(x_i)}\,\eta_{i,\ell}
 \ket{D_N^{(s+2i+\ell h)}},\qquad |\eta_{i,\ell}|=1.
\]
The phases are assigned recursively, as follows.

For $0\le i<k-1$, let $r_{i\ell}=s+2i+\ell h$ and
\[
 L_{i\ell}=\sqrt{w_\ell(x_i)w_\ell(x_{i+1})
 (r_{i\ell}+1)(r_{i\ell}+2)(N-r_{i\ell})(N-r_{i\ell}-1)}.
\]
Take any common diagonal length $d$ in
\[
 \bigl[\max(|L_{i0}-L_{i1}|,|L_{i2}-L_{i3}|),\,
       \min(L_{i0}+L_{i1},L_{i2}+L_{i3})\bigr],
\]
for example its midpoint.  For positive $a,b$ with $|a-b|\le d\le a+b$,
put
\[
 c_a=\frac{d^2+a^2-b^2}{2da},\qquad
 c_b=\frac{d^2+b^2-a^2}{2db},\qquad
 (z_a,z_b)=\bigl(c_a+i\sqrt{1-c_a^2},\,
                   c_b-i\sqrt{1-c_b^2}\bigr).
\]
Then $az_a+bz_b=d$.  Apply this to $(L_{i0},L_{i1})$ to obtain
$(z_0,z_1)$, and to $(L_{i2},L_{i3})$ to obtain $(z_2,z_3)$.
With $\eta_{0,\ell}=1$, define
\[
 (\eta_{i+1,0},\eta_{i+1,1},\eta_{i+1,2},\eta_{i+1,3})
 =(\eta_{i,0}z_0,\eta_{i,1}z_1,-\eta_{i,2}z_2,-\eta_{i,3}z_3).
\]
This is an explicit prescription in radicals, once the interval above is
proved nonempty.

\begin{proof}
Direct substitution gives
\[
 \sum_{\ell=0}^3w_\ell(x)=1,\qquad
 \sum_{\ell=0}^3(x-\ell)w_\ell(x)=0,\qquad
 \sum_{\ell=0}^3(x-\ell)^2w_\ell(x)=v.
\]
The supports are disjoint.  All indices have the same parity, and the gap
between successive bands is $h-2(k-1)=8$.  Consequently all operators
shifting the excitation number by one have zero compression.  For a shift
by two, the only possible off-diagonal entries join rows $i$ and $i+1$
in the same band.  Every such shift has the same orientation across all
four bands.  The phase prescription gives
\[
 \sum_{\ell=0}^3L_{i\ell}\,
 \overline{\eta_{i\ell}}\eta_{i+1,\ell}=0,
\]
so these entries also vanish.  The three moment identities supply
normalization, $PJ_zP=0$, and $PJ_z^2P=j(j+1)P/3$.
Thus all first- and second-rank spin tensors have zero compression,
which is second-order anticoherence of the entire subspace.

It remains to certify that the weights are positive and the quadrilateral
interval is nonempty.  The following exact polynomial certificate does so.
All quantities below are defined by the displayed formulas and require no
numerical approximation.

For each of the eight fixed remainders $\rho$, set
\[
 W_\ell(q,i)=576q^2 w_\ell(x_i).
\]
These are polynomials with integer coefficients.  Under $q=i+z+4$, every
coefficient of each $W_\ell$ is nonnegative, and its constant coefficient
is strictly positive.  This substitution covers the entire basis range
$0\le i\le q-4$, with $z\ge0$, and proves positivity of all weights.

For neighboring rows, define the polynomial
\[
 E_\ell(q,i)=W_\ell(q,i)W_\ell(q,i+1)
 (r_{i\ell}+1)(r_{i\ell}+2)(N-r_{i\ell})(N-r_{i\ell}-1).
\]
It equals $(576q^2)^2L_{i\ell}^2$.
Substitute $q=i+z+5$, which covers all neighboring pairs $0\le i\le q-5$.
Writing $E_\ell$ for these substituted polynomials, put
\[
 F=\left(\sum_\ell E_\ell\right)^2-2\sum_\ell E_\ell^2,
 \qquad A=7421703487488\,(i^{14}z^2+i^2z^{14}).
\]
For each $\rho=0,\ldots,7$, the polynomials
\[
 F+A\quad\hbox{and}\quad64E_0E_1E_2E_3-A^2
\]
have only nonnegative coefficients and strictly positive constant coefficients.  These are finite coefficient assertions;
\texttt{four\_bands.py} constructs the integer polynomials directly and
checks every coefficient, without sampling values of $N$.
Thus $F\ge-A$ and $8\sqrt{E_0E_1E_2E_3}\ge A$ for all $i,z\ge0$, giving
\[
 F+8\sqrt{E_0E_1E_2E_3}>0.
\]
For positive lengths $\lambda_\ell=\sqrt{E_\ell}$, the left side is
\[
 \prod_{\ell=0}^3\left(\sum_{m=0}^3\lambda_m-2\lambda_\ell\right).
\]
At most one factor can be negative.  Therefore this nonnegative product
is equivalent to the four quadrilateral inequalities
$\lambda_\ell\le\sum_{m\ne\ell}\lambda_m$.  Rescaling gives the same
inequalities for $L_{i\ell}$, exactly the condition that the diagonal
interval in the phase prescription has nonempty interior.
\end{proof}

The certificate is recorded in \texttt{four-band-certificate.json}; its
numerical checks of directly constructed spin compressions for
$40\le N\le80$ are supplementary checks of the phase implementation.
The all-$N$ statement follows from the polynomial certificate above.
Combining with the separated-triple construction gives
\[
 K(N,2)\ge\max\!\left\{
 \left\lfloor\frac{N+3}{9}\right\rfloor,
 \left\lfloor\frac N8\right\rfloor-3\right\}\qquad(N\ge40).
\]
\end{document}

% Exact basis calculator: construct_exact.py selects a rational squared diagonal
% by rational bisection and outputs phase definitions without expansion.
